\chapter{Proposed pipeline}
\label{ch:pipeline}

\section{Pipeline overview}
\textit{[Insert one architecture diagram. Show the inference path, the input residual connection, and a separate training-only target and loss branch.]}

\begin{figure}[ht]
\centering
\fbox{\parbox[c][35mm][c]{0.88\textwidth}{\centering
[Pipeline diagram placeholder]\\[3mm]
Show restoration stages, routing weights, input skip,\\
and a separate training-only target and loss path.}}
\caption{Proposed restoration pipeline. Replace this placeholder with the verified architecture diagram.}
\label{fig:pipeline}
\end{figure}

\section{Pipeline components}
\textit{[Describe each component briefly. Cite the source of reused ideas.]}
\begin{singlespace}
\begin{table}[ht]
\centering
\caption{Pipeline component description template.}
\begin{tabular}{p{0.28\textwidth}p{0.62\textwidth}}
\toprule
Component & Details to supply \\
\midrule
Input and degradation & Image range, degradation order, and sampled ranges. \\
Frequency bands & Transform, complementary masks, and cutoffs. \\
Shared experts & Kernel sizes, blocks, channels, and pixel rearrangement. \\
Routing & Gate inputs, top-2 weights, and expert execution. \\
Fusion and output & Channel fusion, refinement, and input residual. \\
\bottomrule
\end{tabular}
\end{table}
\end{singlespace}

\section{Training objective}
\textit{[State the reconstruction, frequency, and routing terms used by the implementation. Define symbols and coefficients. Keep this to one short explanation.]}
